% GENERATED FILE. Edit canonical lessons, then run npm run generate:books.
% canonical-source-hash: fnv1a64:3fdab37d8caffa6d
% canonical-lessons: AR-C339-ghaza, AR-C339-hajama, AR-C339-ihtalla, AR-C339-qada-lead, AR-C339-insahaba

\chapter{Invasions}
\label{ch:ar-invasions}
\hlchaptermodality{339}

\begin{chapteropening}
\textbf{By the end of this chapter:} \emph{I can say to invade, to attack, to occupy, to lead and to withdraw.}

Complete the last lesson of chapter 339: I can say to invade, to attack, to occupy, to lead and to withdraw.
\end{chapteropening}

\section[\texorpdfstring{\ar{غزا} (ghazā)}{غزا (ghazā)}]{\ar{غزا} (ghazā) — to invade, to raid}

\label{lesson:AR-C339-ghaza}



\begin{quote}
Before the new one: say the Arabic for a cemetery, then the Arabic for a coffin.
\end{quote}

\subsection*{You'll want to know: \ar{غزا}}
\textbf{\ar{غزا}} — \emph{ghazā} — \textquotedblleft{}to invade, to raid\textquotedblright{}.

\begin{quote}
\ar{غزا} \ar{الجيش} \ar{البلاد} \ar{المجاورة}.
\end{quote}

\begin{quote}
\emph{ghazā al-jaysh al-bilād al-mujāwira.}
\end{quote}

\begin{quote}
The army invaded the neighbouring lands.
\end{quote}

Present: \ar{يغزو}, \emph{yaghzū}.

Of space: \ar{غزو} \ar{الفضاء}, the conquest of space.

\begin{cousinweb}
\textbf{\ar{غزا}} (\emph{ghazā}) is built on the native Arabic root \emph{gh-z-w} (Lane), \textquotedblleft{}to aim at, to seek\textquotedblright{}, and \textquotedblleft{}to go out to fight\textquotedblright{}; \emph{ghazw} is a raid.
\end{cousinweb}

\begin{grammarlens}[title={What you've built}]
One more word for telling what happened.
\end{grammarlens}

\subsection*{Your turn}
\begin{itemize}
\raggedright
  \item \emph{Say it:} \emph{ghazā}
  \item \emph{Say it:} \emph{ghazā}, once more
  \item \emph{Say it:} \emph{ghazā al-jaysh al-bilād al-mujāwira.} — The army invaded the neighbouring lands.
  \item \emph{Recall:} say the Arabic for a cemetery, then the Arabic for a coffin, then say \emph{ghazā} again
\end{itemize}

\begin{tcolorbox}[breakable,colback=teal!4,colframe=teal!35!black,title={Before you move on}]
What does \ar{غزا} mean? (\textquotedblleft{}To invade, to raid\textquotedblright{}.) Say it once more.
\end{tcolorbox}
\section[\texorpdfstring{\ar{هاجم} (hājama)}{هاجم (hājama)}]{\ar{هاجم} (hājama) — to attack}

\label{lesson:AR-C339-hajama}



\begin{quote}
Before the new one: say the Arabic for a coffin, then the Arabic for to invade.
\end{quote}

\subsection*{You'll want to know: \ar{هاجم}}
\textbf{\ar{هاجم}} — \emph{hājama} — \textquotedblleft{}to attack\textquotedblright{}.

\begin{quote}
\ar{هاجم} \ar{الكلب} \ar{الرجل} \ar{في} \ar{الحديقة}.
\end{quote}

\begin{quote}
\emph{hājama al-kalb ar-rajul fī al-ḥadīqa.}
\end{quote}

\begin{quote}
The dog attacked the man in the garden.
\end{quote}

Present: \ar{يهاجم}, \emph{yuhājim}.

In sport an attacker is \ar{مهاجم}.

\begin{cousinweb}
\textbf{\ar{هاجم}} (\emph{hājama}) is built on the native Arabic root \emph{h-j-m} (Lane), \textquotedblleft{}to come upon someone suddenly\textquotedblright{}. It is Form III, \textquotedblleft{}to rush upon an enemy\textquotedblright{}.
\end{cousinweb}

\begin{grammarlens}[title={What you've built}]
One more word for telling what happened.
\end{grammarlens}

\subsection*{Your turn}
\begin{itemize}
\raggedright
  \item \emph{Say it:} \emph{hājama}
  \item \emph{Say it:} \emph{hājama}, once more
  \item \emph{Say it:} \emph{hājama al-kalb ar-rajul fī al-ḥadīqa.} — The dog attacked the man in the garden.
  \item \emph{Recall:} say the Arabic for a coffin, then the Arabic for to invade, then say \emph{hājama} again
\end{itemize}

\begin{tcolorbox}[breakable,colback=teal!4,colframe=teal!35!black,title={Before you move on}]
What does \ar{هاجم} mean? (\textquotedblleft{}To attack\textquotedblright{}.) Say it once more.
\end{tcolorbox}
\section[\texorpdfstring{\ar{احتل} (iḥtalla)}{احتل (iḥtalla)}]{\ar{احتل} (iḥtalla) — to occupy, to take up}

\label{lesson:AR-C339-ihtalla}



\begin{quote}
Before the new one: say the Arabic for to invade, then the Arabic for to attack.
\end{quote}

\subsection*{You'll want to know: \ar{احتل}}
\textbf{\ar{احتل}} — \emph{iḥtalla} — \textquotedblleft{}to occupy, to take up\textquotedblright{}.

\begin{quote}
\ar{احتل} \ar{فريقنا} \ar{المركز} \ar{الأول}.
\end{quote}

\begin{quote}
\emph{iḥtalla farīqunā al-markaz al-awwal.}
\end{quote}

\begin{quote}
Our team took first place.
\end{quote}

Present: \ar{يحتل}, \emph{yaḥtall}.

Of armies: to occupy a land.

\begin{cousinweb}
\textbf{\ar{احتل}} (\emph{iḥtalla}) is built on the native Arabic root \emph{ḥ-l-l} (Lane), \textquotedblleft{}to untie; to alight in a place\textquotedblright{}. It is Form VIII, \textquotedblleft{}to take up a place\textquotedblright{}.
\end{cousinweb}

\begin{grammarlens}[title={What you've built}]
One more word for telling what happened.
\end{grammarlens}

\subsection*{Your turn}
\begin{itemize}
\raggedright
  \item \emph{Say it:} \emph{iḥtalla}
  \item \emph{Say it:} \emph{iḥtalla}, once more
  \item \emph{Say it:} \emph{iḥtalla farīqunā al-markaz al-awwal.} — Our team took first place.
  \item \emph{Recall:} say the Arabic for to invade, then the Arabic for to attack, then say \emph{iḥtalla} again
\end{itemize}

\begin{tcolorbox}[breakable,colback=teal!4,colframe=teal!35!black,title={Before you move on}]
What does \ar{احتل} mean? (\textquotedblleft{}To occupy, to take up\textquotedblright{}.) Say it once more.
\end{tcolorbox}
\section[\texorpdfstring{\ar{قاد} (qāda)}{قاد (qāda)}]{\ar{قاد} (qāda) — to lead; to drive}

\label{lesson:AR-C339-qada-lead}



\begin{quote}
Before the new one: say the Arabic for to attack, then the Arabic for to occupy.
\end{quote}

\subsection*{You'll want to know: \ar{قاد}}
\textbf{\ar{قاد}} — \emph{qāda} — \textquotedblleft{}to lead; to drive\textquotedblright{}.

\begin{quote}
\ar{قاد} \ar{أبي} \ar{السيارة} \ar{طوال} \ar{الليل}.
\end{quote}

\begin{quote}
\emph{qāda abī as-sayyāra ṭuwāla al-layl.}
\end{quote}

\begin{quote}
My father drove the car all night.
\end{quote}

Present: \ar{يقود}, \emph{yaqūd}.

Of an army: \ar{قاد} \ar{الجيش}, he led the army.

\begin{cousinweb}
\textbf{\ar{قاد}} (\emph{qāda}) is built on the native Arabic root \emph{q-w-d} (Lane), \textquotedblleft{}to lead a horse or camel by its rope\textquotedblright{}; a car is led the same way, so \emph{qāda} is also to drive.
\end{cousinweb}

\begin{grammarlens}[title={What you've built}]
One more word for telling what happened.
\end{grammarlens}

\subsection*{Your turn}
\begin{itemize}
\raggedright
  \item \emph{Say it:} \emph{qāda}
  \item \emph{Say it:} \emph{qāda}, once more
  \item \emph{Say it:} \emph{qāda abī as-sayyāra ṭuwāla al-layl.} — My father drove the car all night.
  \item \emph{Recall:} say the Arabic for to attack, then the Arabic for to occupy, then say \emph{qāda} again
\end{itemize}

\begin{tcolorbox}[breakable,colback=teal!4,colframe=teal!35!black,title={Before you move on}]
What does \ar{قاد} mean? (\textquotedblleft{}To lead; to drive\textquotedblright{}.) Say it once more.
\end{tcolorbox}
\section[\texorpdfstring{\ar{انسحب} (insaḥaba)}{انسحب (insaḥaba)}]{\ar{انسحب} (insaḥaba) — to withdraw, to pull out}

\label{lesson:AR-C339-insahaba}



\begin{quote}
Before the new one: say the Arabic for to occupy, then the Arabic for to lead.
\end{quote}

\subsection*{You'll want to know: \ar{انسحب}}
\textbf{\ar{انسحب}} — \emph{insaḥaba} — \textquotedblleft{}to withdraw, to pull out\textquotedblright{}.

\begin{quote}
\ar{انسحب} \ar{اللاعب} \ar{من} \ar{المباراة}.
\end{quote}

\begin{quote}
\emph{insaḥaba al-lāʿib min al-mubārāh.}
\end{quote}

\begin{quote}
The player pulled out of the match.
\end{quote}

Present: \ar{ينسحب}, \emph{yansaḥib}.

\begin{cousinweb}
\textbf{\ar{انسحب}} (\emph{insaḥaba}) is built on the native Arabic root \emph{s-ḥ-b} (Lane), \textquotedblleft{}to drag along\textquotedblright{}, the root of the taught \emph{saḥaba}. It is Form VII, \textquotedblleft{}to be drawn back\textquotedblright{}.
\end{cousinweb}

\begin{grammarlens}[title={What you've built}]
One more word for telling what happened.
\end{grammarlens}

\subsection*{Your turn}
\begin{itemize}
\raggedright
  \item \emph{Say it:} \emph{insaḥaba}
  \item \emph{Say it:} \emph{insaḥaba}, once more
  \item \emph{Say it:} \emph{insaḥaba al-lāʿib min al-mubārāh.} — The player pulled out of the match.
  \item \emph{Recall:} say the Arabic for to occupy, then the Arabic for to lead, then say \emph{insaḥaba} again
\end{itemize}

\begin{tcolorbox}[breakable,colback=teal!4,colframe=teal!35!black,title={Before you move on}]
What does \ar{انسحب} mean? (\textquotedblleft{}To withdraw, to pull out\textquotedblright{}.) Say it once more.
\end{tcolorbox}
